% Back end: two-level bypass, interlock, execution stages, LSU fast/slow return, long ops, WB.
\begin{tikzpicture}[ares]
% stage headers and pipeline registers
\foreach \n/\c/\x/\w in {ISS/aresTeal/1.45/2.70, EX1/aresAmber/4.35/2.30, EX2/aresAmber/6.95/2.00,
    EX3/aresAmber/9.15/1.60, EX4/aresAmber/10.95/1.20, WB/aresGreen/12.95/2.20}
  \node[stagehead=\c, minimum width=\w cm] at (\x,1.35) {\n};
\foreach \x in {2.99,5.72,8.15,10.15,11.75}
  \node[pbar, minimum height=5.00cm, anchor=north] at (\x,0.45) {};

% ---------------- ISS ----------------
\node[mem=aresTeal, minimum width=1.15cm] (gpr) at (0.70,-0.85) {GPR\\\sub{同步读}};
\draw[fill=white] (1.75,-0.30) -- (2.05,-0.45) -- (2.05,-1.35) -- (1.75,-1.50) -- cycle;
\draw[dat] (gpr.east) -- (1.75,-0.85);
\node[blk=aresTeal, minimum width=2.55cm] (ctl) at (1.45,-2.15)
  {相关检测\\\sub{where·mask → 独热旁路选择}\\\sub{REL \& CONS ≠ 0 → interlock}};
\draw[->, draw=aresTeal!80!black] (1.90,-1.73) -- (1.90,-1.43);
\node[blk=aresTeal, minimum width=1.15cm] (agu) at (0.70,-3.20) {AGU\\\sub{rs1+imm}};

% ---------------- EX1 ----------------
\draw[fill=white] (3.25,-0.30) -- (3.55,-0.45) -- (3.55,-1.35) -- (3.25,-1.50) -- cycle;
\draw[dat] (2.05,-0.90) -- (3.25,-0.90);
\node[blk=aresAmber, minimum width=1.60cm] (fu) at (4.55,-0.90) {ALU · BRU\\CSR · CRC8 前半};
\draw[dat] (3.55,-0.90) -- (fu.west);
\node[blk=aresAmber, minimum width=1.60cm] (md) at (4.55,-4.30) {乘/除法启动};
\node[blk=aresGreen, minimum width=1.60cm] (ls1) at (4.55,-3.20) {LS1\\\sub{PERIP 请求}};
\draw[dat, preaction={draw=white, line width=3pt, -}] (3.40,-1.43) |- (md.west);
\draw[dat] (agu.east) -- node[note, above, pos=0.42]{地址} (ls1.west);

% ---------------- EX2 / EX3 / EX4 ----------------
\node[blk=aresAmber, minimum width=1.85cm] (e2) at (6.95,-0.90) {命中数据选择\\CRC8 后半};
\node[blk=aresAmber, minimum width=1.45cm] (e3) at (9.15,-0.90) {慢 Load\\选择};
\node[blk=aresAmber, minimum width=1.05cm] (e4) at (10.95,-0.90) {传递};
\draw[dat] (fu.east) -- (e2.west);
\draw[dat] (e2.east) -- (e3.west);
\draw[dat] (e3.east) -- (e4.west);
\node[blk=aresGreen, minimum width=1.85cm] (ls2) at (6.95,-3.20) {LS2\\\sub{DCache 命中}};
\node[blk=aresGreen, minimum width=1.45cm] (ls3) at (9.15,-3.20) {LS3\\\sub{桥返回 · 回填}};
\draw[dat] (ls1.east) -- (ls2.west);
\draw[dat] (ls2.east) -- (ls3.west);
\draw[dat] (ls2.north) -- node[note, right, pos=0.45]{快数据} (e2.south);
\draw[dat] (ls3.north) -- node[note, right, pos=0.45]{慢数据} (e3.south);
% long-latency unit
\node[cell=aresAmber, minimum width=5.55cm, minimum height=5mm] (long) at (8.75,-4.30)
  {乘法 4 级部分积流水 · 除法 radix-4 迭代};
\draw[dat] (md.east) -- (long.west);

% ---------------- WB ----------------
\node[blk=aresGreen, minimum width=1.90cm] (sel) at (12.95,-0.90) {结果选择\\\sub{普通 / 乘法 / 除法}};
\node[blk=aresGreen, minimum width=1.90cm] (wr) at (12.95,-2.55) {GPR 写回\\\sub{顺序退休}};
\draw[dat] (e4.east) -- (sel.west);
\draw[dat] (long.east) -- (11.85,-4.30) |- ([yshift=-2.5mm]sel.west);
\draw[dat] (sel.south) -- (wr.north);

% ---------------- bypass bus ----------------
\def\by{0.62}
\draw[fwd, line width=0.9pt, -] (1.90,\by) -- (14.10,\by);
\foreach \x/\y in {5.50/-0.90, 7.95/-0.90, 9.95/-0.90, 11.55/-0.90, 14.10/-0.90}{
  \draw[aresAmber!90!black] (\x,\y) -- (\x,\by);
  \fill[aresAmber!90!black] (\x,\by) circle (1pt);}
\draw[fwd] (1.90,\by) -- (1.90,-0.33);
\draw[fwd] (3.40,\by) -- (3.40,-0.33);
\fill[aresAmber!90!black] (3.40,\by) circle (1pt);
\node[note, text=aresAmber!80!black, anchor=south] at (8.80,\by) {旁路源：EX1 · EX2 · EX3 · EX4 · WB（优先选最年轻的就绪结果）};
\node[note, anchor=west] at (1.96,-0.08) {一级旁路};
\node[note, anchor=west] at (3.46,-0.08) {二级旁路};

% ---------------- interlock note ----------------
\node[note, text=aresRed, anchor=north west, align=left] at (-0.05,-4.80)
  {interlock：保持 IF–ISS，向 EX1 插入气泡；AGU 基址不旁路 EX1，紧邻相关停顿一拍};
\end{tikzpicture}
